%%%%%%%%%%%%%%%%%%%%%%%%%%%%%%%%%%%%%%%%%
% University Assignment Title Page 
% LaTeX Template
% Version 1.0 (27/12/12)
%
% This template has been downloaded from:
% http://www.LaTeXTemplates.com
%
% Original author:
% WikiBooks (http://en.wikibooks.org/wiki/LaTeX/Title_Creation)
%
% License:
% CC BY-NC-SA 3.0 (http://creativecommons.org/licenses/by-nc-sa/3.0/)
%
%%%%%%%%%%%%%%%%%%%%%%%%%%%%%%%%%%%%%%%%%
%\title{Title page with logo}
%----------------------------------------------------------------------------------------
%	PACKAGES AND OTHER DOCUMENT CONFIGURATIONS
%----------------------------------------------------------------------------------------

\documentclass[12pt]{article}
\usepackage[english]{babel}
\usepackage[utf8x]{inputenc}
\usepackage{natbib}
\usepackage{color}
\usepackage[explicit]{titlesec}
\usepackage{url}
\usepackage{graphicx}

\definecolor{dkgreen}{rgb}{0,0.6,0}
\definecolor{gray}{rgb}{0.5,0.5,0.5}
\definecolor{mauve}{rgb}{0.58,0,0.82}

\begin{document}

\begin{titlepage}

\newcommand{\HRule}{\rule{\linewidth}{0.5mm}} % Defines a new command for the horizontal lines, change thickness here

\center % Center everything on the page
 
%----------------------------------------------------------------------------------------
%	HEADING SECTIONS
%----------------------------------------------------------------------------------------

\textsc{\LARGE University of St Andrews}\\[1.5cm] % Name of your university/college
\textsc{\Large CS5014 Coursework 1}\\[0.5cm] % Major heading such as course name
\textsc{\large }\\[0.5cm] % Minor heading such as course title

%----------------------------------------------------------------------------------------
%	TITLE SECTION
%----------------------------------------------------------------------------------------

\HRule \\[0.4cm]
{ \huge \bfseries Machine Learning}\\[0.4cm] % Title of your document
\HRule \\[1.5cm]
 
%----------------------------------------------------------------------------------------
%	AUTHOR SECTION
%----------------------------------------------------------------------------------------


\Large \emph{Author:}\\
 \textsc{150008022}\\[3cm] % Your name

%----------------------------------------------------------------------------------------
%	DATE SECTION
%----------------------------------------------------------------------------------------

{\large \today}\\[2cm] % Date, change the \today to a set date if you want to be precise

%----------------------------------------------------------------------------------------
%	LOGO SECTION
%---------------------------------------------------------------------------------------

\includegraphics[width = 3.1cm]{images/standrewslogo.png}
 
%----------------------------------------------------------------------------------------

\vfill % Fill the rest of the page with whitespace

\end{titlepage}

\part*{Goal}

The goal of this practical was to cleanse and process real world data in order to produce a regression model, and evaluate its performance. \citep{createreactapp}

\part{Loading and Cleaning the Data}

Numpy was used to load the csv file. The header row was skipped, and to ensure there were no missing values, the invalid raise flag was also used when parsing the data. This would raise an exception if any rows were found to be missing data. 

The input and output columns were separated into two variables x and y, as in accordance with the notation used in lectures.

\part{Analysing and Visualising the Data}

Firstly, a histogram of each of the input variables and outputs were plotted in order to visualise the distribution of the values. When comparing to the histograms from the given paper TODO fig 1) ref paper, most of the plots matched. The difference was identified to be caused by 10 bins always being used (the default if not specified by numpy.hist), and the paper would sometimes use more. However, it was still clear that none of the variables had a gaussian distribution.

In order to identify which variables affected the outputs most, scatter graphs were made between each input and output variable. Feature scaling was used to allow for comparisons between values that could have very different ranges. Mean normalisation was used for the scatter plots. The resulting plots resembled those from the given paper (Figure 2 from given paper TODO), with different x-axis values since mean normalisation gave values in the range $-0.5 \leq x \leq 0.5 $.

\part{Choosing Features}

To try and identify which features had the strongest effect on the outputs, both the Pearson and Spearman rank correlation coefficients were used. It was noted that the Pearson correlation would give a perfect value when the two variables were linearly related, whilst the Spearman correlation (a similiar alternative, and the one used in the original paper) would give a perfect value when the variables were monotonically related. From the scatter plots, some variables appeared to have a possible linear relationship with the outputs (for example, X7 and Y1), whilst others had a monotonic relationship (for example, X1 and Y1). Using the Scipy \emph{stats.pearsonr} and \emph{stats.spearmanr} methods, the correlation coefficients were easily calculated, alongside a p-value. 



\part*{Conclusion}

\bibliographystyle{unsrt}
\bibliography{mybib}

\end{document}
